\subsubsection{体位补遗：屈曲位、伸展位与"背入式"的名称说明}

本书前文已系统介绍过主要体位（参见本书体位相关章节），这里补记三个常被提及但在正文中缺少对应条目的名称。

\begin{itemize}
\item \textbf{背入式}：中文资料中"背入式"与"后入式"基本同义——均指女方背对男方、从身后进入的体位。日本 AV 术语体系中的"背面座位""立ちバック"等属同一体位的亚型变体。看到不同名称时不必困惑，它们描述的是同一类动作模式；
\item \textbf{屈曲位}（折り曲げ位）：女方仰卧，双腿向胸部屈曲上抬、膝部贴近肩膀，臀部相应抬起。这一体位使进入角度更接近垂直、插入更深，对阴道前壁与 G 点区域的刺激相对集中；\textbf{风险提示}：插入深度增加意味着对宫颈的顶压更明显，部分女性会感到深部不适或疼痛——出现疼痛应立即调整角度或深度，而不是"忍一忍就过去了"。孕中晚期、宫颈机能不全或有盆底手术史者应避免；
\item \textbf{伸展位}：与屈曲位相对，指身体处于舒展、无屈曲的体位，最典型的是女方仰卧、双腿自然伸展的姿势（如男上位的基本变体）。特点是骨盆角度自然、肌肉放松、易于长时间保持，适合体力有限、腰背不适或需要减轻腹部压迫的人群（如孕中晚期、术后恢复期）——在"舒适度优先"的原则下，它往往比高难度体位更值得推荐。
\end{itemize}

\noindent\textbf{一点原则性提醒}：体位的取舍标准不是"难度"或"新鲜感"，而是\textbf{双方是否舒适、是否有疼痛、是否影响呼吸与循环}。体型差异、关节活动度、既往病史（腰椎、髋关节、心肺疾病）都会影响适配性；"换一个位置"永远比"坚持同一个位置"更专业。
